% -*- coding: utf-8 -*-
% !TEX program = xelatex
\documentclass[plain,noanswer,medmath]{../../template/jnuexam}
\usepackage{../../template/kysx}

\begin{document}


\examtitle{2009}{1}


\makepart{选择题}{1～8小题，每小题4分，共32分}


\begin{problem}
当$x \to 0$时，$f(x) = x - \sin ax$与$g(x) = x^2\ln (1 - bx)$是等价无穷小，则\pickout{}
\begin{abcd}
\item $a = 1,b = - \frac{1}{6}$．
\item $a = 1,b = \frac{1}{6}$．
\item $a = - 1,b = - \frac{1}{6}$．
\item $a = - 1,b = \frac{1}{6}$．
\end{abcd}
\end{problem}


\begin{problem}
如图，正方形$\left\{(x,y) \;\big|\; |x| \le 1, |y| \le 1 \right\}$
被其对角线划分为四个区域$D_k$（$k = 1,2,3,4$），$I_k = \iint_{D_k} y\cos x\dx\dy$，
则$\max\limits_{1 \le k \le 4} \left\{I_k \right\} =$\pickout{}
\par\centerline{\includegraphics[width=0.4\linewidth]{2009A1}}
\begin{abcd}
\item $I_1$．
\item $I_2$．
\item $I_3$．
\item $I_4$．
\end{abcd}
\end{problem}


\begin{problem}
设函数$y = f(x)$在区间$\left[- 1,3 \right]$上的图形为
\par\centerline{\includegraphics[width=0.5\linewidth]{2009A2}}
则函数$F(x) = \int_0^x f(t)\dt$的图形为\pickout{}
\begin{abcd}
\item \includegraphics[width=0.4\linewidth]{2009A3}
\item \includegraphics[width=0.4\linewidth]{2009A4}
\item \includegraphics[width=0.4\linewidth]{2009A5}
\item \includegraphics[width=0.4\linewidth]{2009A6}
\end{abcd}
\end{problem}


\begin{problem}
设有两个数列$\left\{a_n \right\},\left\{b_n \right\}$，若$\lim_{n \to \infty} a_n = 0$，则\pickout{}
\begin{abcd}
\item 当$\sum_{n = 1}^{\infty}b_n$收敛时，$\sum_{n = 1}^{\infty}a_n b_n$收敛．
\item 当$\sum_{n = 1}^{\infty}b_n$发散时，$\sum_{n = 1}^{\infty}a_n b_n$发散．
\item 当$\sum_{n = 1}^{\infty}\left| b_n \right|$收敛时，$\sum_{n = 1}^{\infty}a_n^2b_n^2$收敛．
\item 当$\sum_{n = 1}^{\infty}\left| b_n \right|$发散时，$\sum_{n = 1}^{\infty}a_n^2b_n^2$发散．
\end{abcd}
\end{problem}


\begin{problem}
设$\alpha_1$, $\alpha_2$, $\alpha_3$是$3$维向量空间$R^3$的一组基，
则由基$\alpha_1$, $\frac{1}{2}\alpha_2$, $\frac{1}{3}\alpha_3$到基
$\alpha_1 + \alpha_2$, $\alpha_2 + \alpha_3$, $\alpha_3 + \alpha_1$的过渡矩阵为\pickout{}
\begin{abcd}
\item $\begin{pmatrix}
  1 & 0 & 1 \\
  2 & 2 & 0 \\
  0 & 3 & 3
\end{pmatrix}$．
\item $\begin{pmatrix}
  1 & 2 & 0 \\
  0 & 2 & 3 \\
  1 & 0 & 3
\end{pmatrix}$．
\item $\begin{pmatrix}
   \frac{1}{2} &  \frac{1}{4} & -\frac{1}{6} \\
  -\frac{1}{2} &  \frac{1}{4} & \frac{1}{6} \\
   \frac{1}{2} & -\frac{1}{4} & \frac{1}{6}
\end{pmatrix}$．
\item $\begin{pmatrix}
   \frac{1}{2} & -\frac{1}{2} & \frac{1}{2} \\
   \frac{1}{4} &  \frac{1}{4} & -\frac{1}{4} \\
  -\frac{1}{6} &  \frac{1}{6} & \frac{1}{6}
\end{pmatrix}$．
\end{abcd}
\end{problem}


\begin{problem}
设$A,$$B$均为$2$阶矩阵，$A^*,B^*$分别为$A,$$B$的伴随矩阵，若$|A| = 2,|B| = 3$，
则分块矩阵$\begin{pmatrix}
  O & A \\
  B & O
\end{pmatrix}$的伴随矩阵为\pickout{}
\begin{abcd}
\item $\begin{pmatrix}
     O & 3B^* \\
  2A^* & O
\end{pmatrix}$．
\item $\begin{pmatrix}
     O & 2B^* \\
  3A^* & O
\end{pmatrix}$．
\item $\begin{pmatrix}
     O & 3A^* \\
  2B^* & O
\end{pmatrix}$．
\item $\begin{pmatrix}
     O & 2A^* \\
  3B^* & O
\end{pmatrix}$．
\end{abcd}
\end{problem}


\begin{problem}
设随机变量$X$的分布函数为$F(x) = 0.3\Phi(x) + 0.7\Phi\left(\frac{x - 1}{2}\right)$，
其中$\Phi(x)$为标准正态分布的分布函数，则$EX =$\pickout{}
\begin{abcd}
\item $0$．
\item $0.3$．
\item $0.7$．
\item $1$．
\end{abcd}
\end{problem}


\begin{problem}
设随机变量$X$与$Y$相互独立，且$X$服从标准正态分布$N\left(0,1 \right)$,$Y$的概率分布为
$P\{Y = 0\} = P\{Y = 1\} = \frac{1}{2}$．记$F_Z(z)$为随机变量$Z = XY$的分布函数，
则函数$F_Z(z)$ 的间断点个数为\pickout{}
\begin{abcd}
\item $0$．
\item $1$．
\item $2$．
\item $3$．
\end{abcd}
\end{problem}


\makepart{填空题}{9～14小题，每小题4分，共24分}


\begin{problem}
设函数$f\left(u,v \right)$具有二阶连续偏导数，$z = f(x,xy)$，
则$\frac{\pd^2 z}{\pdx\pd y}=$ \fillin{}．
\end{problem}


\begin{problem}
若二阶常系数线性齐次微分方程$y'' + ay' + by = 0$的通解为$y = (C_1 + C_2 x)\e^x$，
则非齐次方程$y'' + ay' + by = x$满足条件$y(0) = 2,y'(0) = 0$的解为$y =$ \fillin{}．
\end{problem}


\begin{problem}
已知曲线$L:y = x^2\left(0 \le x \le \sqrt{2} \right)$，则$\int_L x \ds =$ \fillin{}．
\end{problem}


\begin{problem}
设$\Omega = \left\{\left. \left(x,y,z \right) \right| x^2 + y^2 + z^2 \le 1 \right\}$，
则$\iiint_{\Omega}z^2 \dx\dy\dz =$ \fillin{}．
\end{problem}


\begin{problem}
若$3$维列向量$\alpha ,\beta$满足$\alpha^T\beta = 2$，
其中$\alpha^T$为$\alpha$的转置，则矩阵$\beta \alpha^T$的非零特征值为 \fillin{}．
\end{problem}


\begin{problem}
设$X_1,X_2, \cdots ,X_m$为来自二项分布总体$B\left(n,p \right)$的简单随机样本，
$\overline{X}$和$S^2$分别为样本均值和样本方差．若$\overline{X} + kS^2$为$np^2$的无偏估计量，
则$k =$ \fillin{}．
\end{problem}


\makepart{解答题}{15～23小题，共94分}


\begin{problem}[points=9]
求二元函数$f(x,y) = x^2\left(2 + y^2 \right) + y\ln y$的极值．
\end{problem}


\begin{problem}[points=9]
设$a_n$为曲线$y = x^n$与$y = x^{n + 1}$（$n = 1,2, \cdots$）所围成区域的面积，
记$S_1 = \sum_{n = 1}^{\infty}a_n$，$S_2 = \sum_{n = 1}^{\infty}a_{2n - 1}$，
求$S_1$与$S_2$的值．
\end{problem}


\begin{problem}[points=11]
椭球面$S_1$是椭圆$\frac{x^2}{4} + \frac{y^2}{3} = 1$绕$x$轴旋转而成，
圆锥面$S_2$是由过点$\left(4,0 \right)$且与椭圆$\frac{x^2}{4} + \frac{y^2}{3} = 1$
相切的直线绕$x$轴旋转而成．\par
\begin{enumerate*}
\item 求$S_1$及$S_2$的方程；
\item 求$S_1$与$S_2$之间的立体体积．
\end{enumerate*}
\end{problem}


\begin{problem}[points=11]
\begin{enumerate}
\item 证明拉格朗日中值定理：若函数$f(x)$在$\left[a,b \right]$上连续，在$(a,b)$可导，
      则存在$\xi \in \left(a,b \right)$，使得$f(b) - f(a) = f'\left(\xi \right)(b - a)$．
\item 证明：若函数$f(x)$在$x = 0$处连续，在$\left(0,\delta \right)\left(\delta > 0 \right)$内可导，
      且$\lim_{x \to 0^+} f'(x) = A$，则$f'_+(0)$存在，且$f'_+(0) = A$．
\end{enumerate}
\end{problem}


\begin{problem}[points=10]
计算曲面积分$I = \oiint_{\Sigma}\frac{x\dy\dz + y\dz\dx + z\dx\dy}{\left(x^2+y^2+z^2\right)^{\frac{3}{2}}}$，
其中$\Sigma$是曲面$2x^2 + 2y^2 + z^2 = 4$的外侧．
\end{problem}


\begin{problem}[points=11]
设$A = \begin{pmatrix}
   1 & -1 & -1 \\
  -1 &  1 & 1 \\
   0 & -4 & -2
\end{pmatrix}$,$\xi_1 = \begin{pmatrix}
  -1 \\
  1 \\
  -2
\end{pmatrix}$．
\begin{enumerate}
\item 求满足$A\xi_2 = \xi_1,A^2\xi_3 = \xi_1$的所有向量$\xi_2,\xi_3$；
\item 对(I)中的任意向量$\xi_2,\xi_3$，证明$\xi_1,\xi_2,\xi_3$线性无关．
\end{enumerate}
\end{problem}


\begin{problem}[points=11]
设二次型$f\left(x_1,x_2,x_3 \right) = ax_1^2 + ax_2^2 + (a - 1)x_3^2 + 2x_1x_3 - 2x_2x_3$．
\begin{enumerate}
\item 求二次型$f$的矩阵的所有特征值；
\item 若二次型$f$的规范形为$y_1^2 + y_2^2$，求$a$的值．
\end{enumerate}
\end{problem}


\begin{problem}[points=11]
袋中有$1$个红球，$2$个黑球与$3$个白球．现有放回地从袋中取两次，每次取一个球，
以$X,Y,Z$分别表示两次取球所取得的红球、黑球与白球的个数．\par
\begin{enumerate*}
\item 求$P\left\{X = 1 | Z = 0 \right\}$；
\item 求二维随机变量$(X,Y)$的概率分布．
\end{enumerate*}
\end{problem}


\begin{problem}[points=11]
设总体$X$的概率密度为$f(x) = \begin{cases}
  \lambda^2 x\e^{- \lambda x}, & x > 0, \\
                            0, & \text{其他},
\end{cases}$其中参数$\lambda$（$\lambda > 0$）未知，
$X_1$, $X_2$, $\cdots$, $X_n$是来自总体$X$的简单随机样本．\par
\begin{enumerate*}
\item 求参数$\lambda$的矩估计量；
\item 求参数$\lambda$的最大似然估计量．
\end{enumerate*}
\end{problem}


\end{document}
